%!TEX root = main.tex
%=========================================================

\section{Interoperability mechanisms}

\label{sec:mechanisms}
\begin{itemize}
	\item we follow a main switching criterion: nodes in dense areas maintain a dissemination overlay and nodes in sparse areas broadcast via a controlled flooding technique \er{this should have been said in the previous section, not here}
	\item our adaptive approach is safe, i. e., broadcast messages are delivered to every node that is part of the wireless topology \raziel{we cannot say to every node at the communication area because connectivity between nodes is not a guarantee due to mobility}
	\item in other words, our approach is as good as running one single broadcast technique in the whole network \er{I would not say it this way: Our approach ensures continuity in the delivery despite the fact that some regions use different protocols}
	\item our adaptive approach is safe because we will prove that broadcast messages are delivered to every node in the network when source nodes: \er{this will need a figure to illustrate the point}
	\begin{enumerate}
		\item are part of a dissemination overlay, or
		\item use controlled flooding
	\end{enumerate}
	\item prove 1 and 2 is equal to prove that messages are correctly forwarded when messages transit between any pair of nodes running different protocols \er{we got rid of this way of thinking and wanted to discuss about regions and going from a region to the other. This 'general interoperability' claim is too bold and too general.}
	\begin{itemize}
		\item once this transition takes place, delivery of messages among nodes running the same protocol is guaranteed by the safety properties of the protocol itself
	\end{itemize}
	\item let be $O$, $F$ nodes where $O$ is part of a dissemination overlay and $F$ runs a controlled flooding approach \er{inconsistent used of O and F, are these sets, nodes or protocols? No need to formalize now.}
	\begin{itemize}
		\item overlay-based techniques tagged $O$ as a node being part or not of the forwarding set \textit{FwdSet}
		\item $F$ forwards incoming messages iif its forwarding condition isn't reached yet
	\end{itemize}
	\item transition of message $F \rightarrow O$
	\begin{itemize}
		\item \textbf{$O \in FwdSet$ ($O$ is a relay within the dissemination overlay):}  incoming messages, that haven't been previously heard are forwarded following the overlay approach. \textbf{OK1}
		\item \textbf{$O \notin FwdSet$ ($O$ is not a relay):} here we tagged $O$ as a interoperable relay (border node?) \er{use border node}. The message is transformed in an overlay-based packet but forwarded using a simple flooding approach. \textbf{OK2} \raziel{this is done to avoid changing the structure of the overlay}. Or may be forwarded as if $O$ is a new source. \raziel{this may be possible but will have a bad impact in the prove} \er{I do not think you need the latter}
	\end{itemize}
	\item transition of message $O \rightarrow F$
	\begin{itemize}
		\item \textbf{$O \in FwdSet$ ($O$ is a relay within the dissemination overlay):} here again $F$ receive the message but forward it following the flooding approach \textbf{OK3}
	\end{itemize}
	\item still we need to prove two cases, a proof for one case is valid for the other\raziel{I do not follow the previous format because these cases are very similar}
		\begin{itemize}
			\item remind again that $F$ and $O$ are in range 
		\end{itemize}
		\begin{enumerate}
			\item $F$ reaches its forwarding condition, then, the message will not be delivered to $O$
			\item $O$ is not a relay in the overlay, then, the message will not be delivered to $F$
		\end{enumerate}
	\item here, we need to remember that an adaptation policy requires the exchange of control messages to apply such policy otherwise we will never have a situation where different protocols are running in the network, i. e, the switching criterion already took place. \raziel{this require to present first the adaptation approach and then the problems of interoperability cause exchange of beacons is required to spot interoperable relays}
	\er{this seems overly complicated to me. What you want to say first is that you assume that border nodes know they are border nodes, and simply explain why your protocol works when it is the case (no need for excessive formalism). Then you should justify that the key problem is for nodes to detect when they are border nodes and avoid too many false positives for performance reasons.}
	\item This means that the closure coefficient \er{not sure what it has to do here} was previously computed or at least $F$ and $O$ have performed an exchange of control messages in the past. We called such exchange a history $H$
		\item \er{the policy should go in the following section not here. You should only say that when the selection of border nodes is correct then the interoperability works.}	
		we propose an interoperability policy based on nodes degree to deal with 1 and 2: if the local closure coefficient detects a border node, i. e. closure coefficient is in the threshold of medium density, and $H$ is not empty then, $F$ or $O$ turn to interoperable relays. \textbf{OK4} \raziel{of course, here we can show the distribution of the local closure coefficient of nodes located at the border to give a sort of proof in form of analytic study}
		\item due to OK$[1-4]$ our approach is safe
\end{itemize}

\er{
The goal of this section is to present the technical contributions that allow switching from one protocol to the other:
\begin{itemize}
	\item The mechanism by which a system with different protocols still guarantee the dissemination:
	\begin{itemize}
		\item Showing examples of why the straightforward approach does not work (I would suggest a reasoning based on system physical boundaries: the overlay-based approach considers that it spreads the message within the boundaries of the set of nodes that employ this mechanisms, and therefore ignores that there are exterior nodes that should also receive the message)
		\item Introduce the role of interoperability relay nodes (or find another easy-to-understand name) and that we will need to identify at runtime these nodes and assign them with the role of forwarding from one system/protocol to another system/protocol. We note that this role is necessary only from the overlay-based to the flooding-based, but not in the other direction. Another way to present it is to say that the emergent overlay has to ``connect'' to the surrounding flooding-based system, hence needs connection points.\raziel{interoperability should be address in both directions, for instance, nodes that aren't relays in dense zones may drop incoming messages from others running a controlled flooding approach.}
		\item Showing the algorithm used to detect and enable interoperability connection points.
		\item Prove, at least informally, that the algorithm is correct.\raziel{we do not have an algorithm per se, the technique that shows better results to spot interoperability relay nodes is when they do have a medium density based on their degree; we use the local closure coefficient as measure of degree}
		\item Discuss its potential limitations, in particular when the overlay-based zones are very small, in which case there will be more nodes that will have to act as connection points than nodes that are relays, which defeats the purpose
		\item Discuss how connection points evolve when the system is subject to mobility
	\end{itemize}
	\item The grace period mechanism that ensures that the switch between flooding and overlay-based broadcast does not lead to missed messages:
	\begin{itemize}
		\item Using an example, show that the immediate switch, before the overlay creation protocol includes the node in the overlay, can lead to missed messages
		\item On the same example, show that waiting for a bit of time, or under a certain condition is met (received at least one beacon from the overlay creation? Or do we need something more complex?), while participating to the two protocols, avoids the problem
	\end{itemize}
\end{itemize}
}
